% !PRINTED 539
\weberpage{539}
%<<CONT
\begin{enumerate}
\item[] und umgekehrt:
\item[2.] \gesperrt{Ist $\beta_1, \beta_2, \ldots, \beta_n$ ein System ganzer durch $\mu$
theilbarer Zahlen, dessen Discriminante gleich
ist dem Quadrate der absoluten Norm von $\mu$,
multiplicirt mit der Grundzahl des Körpers, so
ist $\beta_1, \beta_2, \ldots, \beta_n$ eine Basis von $\mu$.}
\end{enumerate}

\S.\,148.

\gesperrt{Volles Restsystem nach einem Modul.}

\begin{enumerate}
\item[1.] \textbf{Definition:} \gesperrt{Zwei ganze algebraische Zahlen
$\xi, \eta$, deren Differenz $\xi - \eta$ durch ein ganzes
von Null verschiedenes Functional $\mu$ theilbar
ist, heissen mit einander congruent nach dem
Modul $\mu$.}
\end{enumerate}

Der Begriff der Congruenz lässt sich auch auf Functionale
ausdehnen, was wir aber fürs Erste noch nicht thun. Dagegen
ist es wesentlich, als Moduln der Congruenz nicht bloss die
Zahlen, sondern auch die Functionale zu berücksichtigen. Jede
Congruenz bleibt bestehen, wenn der Modul durch ein associirtes
Functional ersetzt wird. Beim Modul können beliebige Einheits\-
factoren hinzugefügt oder weggelassen werden.

Wir gebrauchen für die Congruenz das \gesperrt{Gauss}'sche Zeichen
(Bd. I, \S.\,115)
\weq{1}{\xi \equiv \eta \ (\mathrm{mod}\ \mu).}

Eine solche Congruenz ist gleichbedeutend mit der Gleichung
\weq{2}{\xi = \eta + \omega\mu,}
worin $\omega$ irgend ein ganzes Functional sein kann.

Aus dieser Darstellung erkennt man dann sofort, dass, ebenso
wie in Congruenzen zwischen rationalen Zahlen, wenn man in
einem durch \gesperrt{Addition}, \gesperrt{Subtraction} und \gesperrt{Multiplication} von
ganzen Zahlen zusammengesetzten Ausdruck jede Zahl durch
eine nach dem Modul $\mu$ congruente Zahl ersetzt, eine nach dem\-
selben Modul congruente Zahl das Resultat ist; in Zeichen:

Sind $\xi_1, \eta_1, \xi_2, \eta_2, \ldots$ ganze Zahlen, die den Congruenzen
\[
\xi_1 \equiv \eta_1, \ \xi_2 \equiv \eta_2, \ldots \ (\mathrm{mod}\ \mu)
\]
%CONT>>
